% Q09 -- 3T DRAM cell schematic.
% M1 = write-access nMOS (gate=WL, terminal driven by write bit line BLbar) charges node X;
% X = gate of M2 (read-evaluation nMOS); M2 in series with M3 = read-access nMOS (gate=RL);
% read bit line BL at the top of the series branch.
\documentclass[border=10pt]{standalone}
\usepackage{circuitikz}
\begin{document}
\begin{circuitikz}
% ---- write-access transistor M1 ----
\path (-3,3) node[nmos] (M1){};
\node[right] at (M1.D) {M1};
\draw (M1.D) -- (-3,5.2) node[above]{$\overline{BL}$ (write)};
\draw (M1.G) -- (-5.2,3) node[left]{WL};
% ---- storage node X (M1 source -> M2 gate) ----
\draw (M1.S) -- (-3,1.3) coordinate (xa);
\node[circ] at (xa) {};
\node[below] at (-3,1.2) {node X};
\draw (xa) -- (0.3,1.3);                 % run across to M2 gate level
% parasitic storage capacitance at X
\draw (-1.5,1.3) to[C, l=$C_S$] (-1.5,-0.6);
\draw (-1.5,-0.6) node[ground]{};
% ---- read-evaluation transistor M2 (gate = X) ----
\path (2,1.5) node[nmos] (M2){};
\node[right] at (M2.D) {M2};
\draw (M2.G) -| (0.3,1.3);                % gate tied to node X
\draw (M2.S) -- (2,-0.4) node[ground]{};
% ---- read-access transistor M3 (gate = RL), in series above M2 ----
\path (2,4.3) node[nmos] (M3){};
\node[right] at (M3.D) {M3};
\draw (M3.S) -- (M2.D);                   % series connection
\draw (M3.G) -- (-0.2,4.3) node[left]{RL};
\draw (M3.D) -- (2,6.2) node[above]{BL (read, precharged)};
\end{circuitikz}
\end{document}
